\begin{table}[t]
\centering
\caption{Isolated actor costs and initialization savings}
\label{tab:r5_actor_timing}
\small
\begin{tabular}{rrrrrr}
\toprule
$d$ & Seed & Actor train & Actor query & Zero query & Projected break-even \\
\midrule
4 & 101 & 1.483 & 0.0414 & 0.0424 & 48512 \\
4 & 202 & 0.741 & 0.0432 & 0.0451 & 12544 \\
4 & 303 & 0.742 & 0.0436 & 0.0429 & None \\
8 & 101 & 0.808 & 0.0607 & 0.0626 & 13568 \\
8 & 202 & 0.797 & 0.0637 & 0.0713 & 3392 \\
8 & 303 & 0.794 & 0.0489 & 0.0512 & 11488 \\
16 & 101 & 0.878 & 0.1004 & 0.1035 & 8864 \\
16 & 202 & 0.879 & 0.1174 & 0.1326 & 1856 \\
16 & 303 & 0.875 & 0.1051 & 0.1079 & 9888 \\
\bottomrule
\end{tabular}
\par\medskip
\begin{minipage}{.97\linewidth}
\footnotesize Seconds per 32-state complete-tree query batch. Actor training is the distillation fit alone; common target construction and critic training are not charged twice. The last column is $32\lceil T_{\rm actor}/(t_{\rm zero}-t_{\rm actor})\rceil$ when the measured saving is positive. It is a linear fixed-batch projection, not an observed crossover. ``None means that actor initialization was not faster in that timing comparison. Optimizer iteration counts and separated initialization times are saved in the machine-readable results.
\end{minipage}
\end{table}
